%! Summary Statistics for a Quantitative Variable — Unit 1, Lesson 1.5 — Warm-Up (Answer Key)
% ONE PAGE, blank and key. Three items.
% Mirrors the blank exactly apart from the package swap, the pageheader, and the answers.
% NO teachernote here — teacher prose lives in the lesson plan.
\documentclass[10pt]{article}
\usepackage{apstats-article}
\usepackage{apstats-key}   % pulls in -boxes; do NOT also load -boxes

\begin{document}

\pageheader{Unit 1, Lesson 1.5}{Warm-Up --- Answer Key}

\vspace{0.10in}

\begin{enumerate}[leftmargin=1.7em, itemsep=0.17in]

  \item \textbf{From last lesson.} Twelve students reported how many hours of television they
        watched last week.

        \vspace{2pt}
        \begin{center}
        \begin{tikzpicture}[scale=0.95]
          \draw[->, thick] (-0.2,0) -- (7.1,0);
          \foreach \x/\lab in {0/0, 0.9/2, 1.8/4, 2.7/6, 3.6/8, 4.5/10, 5.4/12, 6.3/14} {
            \draw (\x,-0.07) -- (\x,0.07); \node[below, font=\scriptsize] at (\x,-0.1) {\lab};}
          \foreach \x/\n in {0.45/2, 0.9/3, 1.35/2, 1.8/2, 2.25/1, 2.7/1, 6.3/1} {
            \foreach \k in {1,...,\n} { \fill[navy] (\x,{0.2*\k}) circle (0.075); }}
          \node[font=\scriptsize] at (3.4,-0.55) {Hours of television};
        \end{tikzpicture}
        \end{center}
        \vspace{2pt}

        Describe the \textbf{shape} of this distribution, and name any value that stands apart
        from the rest.\\[4pt]
        \ansline{Unimodal and skewed to the right: most students watched 1--4 hours and the}\\[1pt]
        \ansline{tail stretches to the high end. The student with 14 hours stands well apart.}

  \item Here are six numbers: \quad 4 \quad 9 \quad 6 \quad 12 \quad 7 \quad 10

        \vspace{5pt}
        Their \textbf{average}: \ans{8} \qquad
        Put them in order and give the value in the \textbf{middle}: \ans{8 (between 7 and 9)}

  \item Here are seven numbers, already in order.

        \vspace{4pt}
        \begin{center}
        3 \qquad 5 \qquad 8 \qquad 9 \qquad 12 \qquad 15 \qquad 20
        \end{center}
        \vspace{4pt}

        \textbf{(i)} Which one is in the middle? \ans{9} \quad
        How many are below it? \ans{3} \quad Above it? \ans{3}
        \vspace{5pt}

        \textbf{(ii)} Now look only at the numbers \emph{below} the middle one. Which of
        \emph{those} is in the middle? \ans{5}
        \vspace{5pt}

        \textbf{(iii)} And looking only at the numbers \emph{above} the middle one, which of
        those is in the middle? \ans{15}

\end{enumerate}

\end{document}
